% Section 7 — Conclusion
% Data source: Phase 2 SciERC registry experiments + Phase 0 DocRED
% exp_ids: phase2_scierc_unconstrained, phase2_scierc_full_s1,
%          phase0_scierc_unconstrained, phase0_docred_unconstrained
% All numbers verified against registry 2026-05-14.

We presented OntoCoDec, a framework that automatically compiles IE ontology schemas into token-level decoding constraints to enforce type validity at generation time.
Our work makes three contributions.

First, we conducted the first systematic study of ontological constraint violations in LLM-based IE (\S\ref{sec:violation_landscape}).
Across two benchmarks with substantially different ontology sizes---SciERC (7~relations) and DocRED (96~relations)---we found that domain/range violation rates are comparable (15.37\% vs.\ 18.46\%), % exp_ids: phase0_scierc_unconstrained, phase0_docred_unconstrained
indicating that violation severity is driven by \emph{constraint tightness} $\bar{\tau}$ rather than ontology scale.
SciERC ($\bar{\tau}{=}0.333$, mean 24 valid type pairs per relation) and DocRED ($\bar{\tau}{=}0.842$, mean 5.68) exhibit comparable violation rates despite a 14$\times$ difference in ontology size, and within each model, DR reduction scales inversely with $\bar{\tau}$: 94--99\% on SciERC ($\bar{\tau}{=}0.333$) vs.\ 45--90\% on DocRED ($\bar{\tau}{=}0.842$), with the residual modulated by model-specific baseline violation rates (Section~\ref{sec:cross-model}).
% Qwen: SciERC DR↓94%, DocRED DR↓45%; Llama: SciERC DR↓99%, DocRED DR↓90%

Second, we introduced a two-layer constraint architecture that decomposes ontological enforcement into complementary mechanisms:
a \emph{grammar layer} (XGrammar CFG) that eliminates structural violations (RV, ET\,$\to$\,0\%) and a \emph{semantic layer} (TypeConstraintProcessor with a stateful entity registry) that enforces domain/range constraints at the token level. % exp_id: phase2_scierc_grammar_only
The grammar layer serves as an enabling prerequisite for the semantic layer; the two compose via bitwise AND over candidate token sets.

Third, we validated OntoCoDec empirically.
On SciERC, the full system reduces domain/range violations from 15.37\% to 0.93\%---a 94\% relative reduction---while completely eliminating RV and ET violations, with no significant impact on relation F1 ($+$0.13 percentage points, $p > 0.05$). % exp_id: phase2_scierc_full_s1
On DocRED, the full system reduces domain/range violations from 18.46\% to 10.21\%---a 45\% relative reduction---while eliminating RV and ET violations entirely, with statistically significant gains in relation F1 ($+$0.50pp, $p{<}0.001$) and entity F1 ($+$1.23pp, $p{<}0.001$) for Qwen3-8B. Llama-3.1-8B achieves substantial DR reduction (90.2\%, 26.07\%\,$\to$\,2.55\%) with significant R-F1 gains ($+$0.29pp, $p{=}0.013$) and a modest E-F1 trade-off ($-$0.64pp, $p{=}0.022$). % exp_ids: phase2_docred_full, llama_docred_full_s1, llama_docred_unconstrained
Compared with post-hoc filtering, decoding-time enforcement preserves true positives and eliminates entity-type violations that post-hoc methods cannot address.

\paragraph{Limitations.}\label{sec:limitations}
Several limitations qualify these findings.
The semantic layer introduces cross-field consistency (CF) violations: CF rises from 0.08\% under unconstrained decoding to 9.90\% with full constraints. % exp_ids: phase2_scierc_unconstrained, phase2_scierc_full_s1
This is a systematic consequence of per-relation independent type enforcement---when the same entity appears across multiple triples, each relation's domain/range constraint may assign it a different type, producing entity-versus-relation type drift.
The grammar layer alone does not elevate CF (0.08\%, matching the unconstrained baseline), confirming that CF is entirely attributable to the semantic layer. % exp_id: phase2_scierc_grammar_only
Resolving this cross-relation type coherence problem is an open challenge, not a trivial implementation fix.
Cross-field consistency violations increase with constraint strength and vary across models: CF reaches 12.02\% for Llama-3.1-8B on DocRED under full constraints, compared with 2.57\% for Qwen3-8B on the same dataset, suggesting that per-relation independent enforcement interacts with model-specific generation patterns. % exp_ids: llama_docred_full_s1, phase2_docred_full
Constrained decoding can increase parse failure rates for some model-dataset combinations: Llama-3.1-8B on DocRED shows 91 parse failures under full constraints versus 68 unconstrained (9.1\% vs.\ 6.8\%), likely due to max-token truncation of longer constrained outputs. All F1 metrics are computed over successfully parsed documents only. % exp_ids: llama_docred_full_s1, llama_docred_unconstrained
Beyond CF, experiments use two models (Qwen3-8B and Llama-3.1-8B-Instruct), confirming that constraint enforcement generalizes across model families: Llama-3.1-8B on DocRED achieves a 90.2\% DR reduction (26.07\%\,$\to$\,2.55\%) without any code modification.
% llama_docred_unconstrained: DR=26.07%; llama_docred_full_s1: DR=2.55%
Results are reported for two models and two datasets under zero-shot prompting only, and integration with fine-tuned models has not been explored.

\paragraph{Future Work.}
We identify four directions.
(1)~Extension to larger models (quantized 70B-class) to further verify model-agnostic generalization across scales.
(2)~Global entity-level type consensus during decoding---maintaining a shared entity registry across all relations within an instance---to resolve CF violations.
(3)~Integration with fine-tuned models, combining constraint-aware training objectives with decoding-time enforcement.
(4)~Extension to non-English ontologies to evaluate cross-lingual applicability.
